%=============================================================================
% Wave 4, Iteration 18: P12 (EM-Psi Polariton) + P15 (Detection)
% Agent 13: P12 niche + P15 ranking crisis triage
%
% Model: Opus 4.6
% Date: 20 March 2026
%
% INHERITED STATE (from i17):
%   - P12 REVIVED from DEAD to NICHE (coherence length 0.91m at 47 GHz)
%   - Multi-cavity superradiance N=3, C_3=3.44
%   - P15 detection: Process A unaffected by superconductivity
%   - Process B suppressed by BCS gap
%   - Optimal phase gate |alpha/beta|* = 1+sqrt(2)
%   - EXISTENTIAL: Dust-branch c_s ~ c/2.
%
% MANDATE (i18):
%   P12 niche survives? P15 detection ranking updated?
%=============================================================================

\documentclass[11pt]{article}
\usepackage{amsmath,amssymb}
\usepackage{geometry}
\geometry{margin=1in}
\usepackage{booktabs}
\usepackage{xcolor}
\usepackage{tcolorbox}

\title{\textbf{Wave 4, Iteration 18:}\\
\textbf{P12 + P15 --- Crisis Triage}\\[6pt]
{\large P12 Niche Survives; P15 Lab Detection Immune}}
\author{DFD Research Programme --- W4i18 Agent 13 (Opus 4.6)}
\date{20 March 2026}

\begin{document}
\maketitle

%=============================================================================
\section*{Executive Summary}
%=============================================================================

\begin{tcolorbox}[colback=green!5!white, colframe=green!70!black,
  title=\textbf{W4i18: P12 NICHE SURVIVES, P15 LAB DETECTION IMMUNE}]

\begin{itemize}
\item \textbf{SURVIVES (P12)}: EM-psi polariton at 47\,GHz, single-cavity
  niche (0.91\,m coherence), multi-cavity superradiance $N=3$. All
  derive from local $\psi$--EM coupling.
\item \textbf{SURVIVES (P15)}: Process A (EM$\to\psi$ sourcing),
  optimal phase gate, laboratory detection protocols. All lab physics.
\item \textbf{AT RISK}: P12's atmospheric/space propagation applications
  (already dead). P15's cosmological detection channels (bright sirens,
  CMB lensing) depend on cosmological $\psi$ predictions.
\item \textbf{Crisis contribution}: P12's polariton measurement would
  confirm EM--$\psi$ coupling strength directly.
\item \textbf{TOP FINDING}: Both P12 (niche) and P15 (lab detection)
  are crisis-immune in their surviving forms. P12's revival from DEAD
  to NICHE in i17 is unaffected by the dust-branch crisis.
\end{itemize}
\end{tcolorbox}

%=============================================================================
\section{P12: EM-Psi Polariton --- Niche Survives}
%=============================================================================

P12's niche application (single-cavity polariton) depends on:
\begin{itemize}
\item Avoided crossing at $\sim 47$\,GHz between EM and $\psi$ modes.
\item Rabi splitting $\Delta f = 1.83$\,GHz.
\item Coherence length 0.91\,m (restored in i17 after $c_\psi = c$
  confirmation).
\end{itemize}

These are all \textbf{local electromagnetic--$\psi$ coupling effects}
measured in a laboratory cavity. The dust-branch crisis concerns
cosmological perturbation growth, which has no bearing on cavity
polariton physics.

\begin{table}[h]
\centering
\begin{tabular}{lcc}
\toprule
\textbf{P12 Prediction} & \textbf{Cosmo?} & \textbf{Status} \\
\midrule
Avoided crossing 47\,GHz & No & SURVIVES \\
$\Delta f = 1.83$\,GHz Rabi & No & SURVIVES \\
Coherence 0.91\,m & No & SURVIVES \\
$N=3$ superradiance $C_3 = 3.44$ & No & SURVIVES \\
Atmospheric propagation & Dead (pre-crisis) & DEAD \\
Space-to-space WPT & Dead (pre-crisis) & DEAD \\
\bottomrule
\end{tabular}
\end{table}

%=============================================================================
\section{P15: Detection Programme --- Lab Channels Immune}
%=============================================================================

P15 detection channels split into:

\subsection{Crisis-Immune (Lab) Channels}
\begin{itemize}
\item \textbf{Process A} (EM$\to\psi$ sourcing): Direct laboratory
  excitation of $\psi$-modes. Independent of cosmology.
\item \textbf{Optimal phase gate}: $|\alpha/\beta|_* = 1 + \sqrt{2}$.
  Algebraic result, no cosmological input.
\item \textbf{SQMS cavity measurements}: Q-ratio, $\psi$-modulation.
\item \textbf{MEMS birefringence}: Local $\psi$-field effect.
\item \textbf{TiN kinetic inductance}: Material property measurement.
\end{itemize}

\subsection{Cosmology-Dependent Channels (At Risk)}
\begin{itemize}
\item \textbf{Bright siren $H_0$}: Requires DFD prediction for
  luminosity distance $d_L(z)$, which depends on cosmological $\psi$.
\item \textbf{CMB lensing $A_L$ split}: Requires cosmological
  $\psi$-field lensing contribution. AT RISK.
\item \textbf{$S_8$ tomographic gradient}: Requires cosmological
  $P(k)$ prediction. AT RISK.
\item \textbf{BAO shift}: Requires dust-branch modification. AT RISK.
\end{itemize}

\subsection{Updated Detection Ranking (Post-Crisis)}

\begin{enumerate}
\item \textbf{SQMS Phase 0} (Q-ratio, \$47--50K, 2 weeks) --- IMMUNE
\item \textbf{Cat-psi protocol} (\$200--400K, 6 months) --- IMMUNE
\item \textbf{LLR at APOLLO} (0.93\,ns, existing hardware) --- IMMUNE
\item \textbf{MEMS birefringence / TiN readout} --- IMMUNE
\item \textbf{P7 Phase 0 energy storage} (\$1.55M) --- IMMUNE
\item \textbf{Einstein Telescope GW--$\psi$} ($\sim$2035+) --- IMMUNE
\item \textbf{Euclid $S_8$ tomography} (Oct 2026) --- AT RISK$^*$
\item \textbf{DESI BAO} --- AT RISK$^*$
\item \textbf{Bright sirens (ET+CE)} --- AT RISK$^*$
\end{enumerate}

$^*$These become moot or need revised predictions if \texttt{mochi\_class}
shows cosmological sector failure.

%=============================================================================
\section*{TOP FINDING}
%=============================================================================

\begin{tcolorbox}[colback=blue!5!white, colframe=blue!70!black]
\textbf{P12's niche revival (i17) is crisis-immune: polariton physics
is lab-local. P15's detection programme retains 6 of 9 channels as
crisis-immune (all lab-based). The top 6 detection experiments proceed
regardless of cosmological outcome. The bottom 3 (Euclid, DESI, bright
sirens) await \texttt{mochi\_class} resolution.}
\end{tcolorbox}

\end{document}
